\documentclass[11pt,a4paper]{article}
\usepackage{fontspec,amsmath,amssymb,amsthm,microtype}
\setmainfont{lmroman10-regular.otf}[BoldFont=lmroman10-bold.otf,ItalicFont=lmroman10-italic.otf,BoldItalicFont=lmroman10-bolditalic.otf]
\usepackage{xcolor}
\usepackage[most]{tcolorbox}
\usepackage{mathrsfs}
\providecommand{\lt}{<}
\providecommand{\gt}{>}
\setmonofont{DejaVuSansMono.ttf}[BoldFont=DejaVuSansMono-Bold.ttf,ItalicFont=DejaVuSansMono-Oblique.ttf,BoldItalicFont=DejaVuSansMono-BoldOblique.ttf,Scale=0.86]
\newfontfamily\leanfallback{STIXTwoMath-Regular.otf}[Scale=0.86]
\definecolor{leanrule}{gray}{0.62}
\definecolor{leanback}{gray}{0.962}
\definecolor{leantitle}{gray}{0.22}
\definecolor{leanmuted}{gray}{0.38}
\newtcolorbox{leanbox}[1]{enhanced,breakable,colback=leanback,colframe=leanrule,
  boxrule=0.4pt,arc=0.6mm,left=2.4mm,right=2.4mm,top=1.8mm,bottom=1.8mm,
  before skip=0.8em,after skip=1em,fonttitle=\small\bfseries,coltitle=leantitle,
  colbacktitle=leanback,titlerule=0.3pt,title={#1}}
\newlength{\leancharwidth}
\newenvironment{leancode}{\par\smallskip\begingroup\ttfamily\footnotesize\raggedright
  \settowidth{\leancharwidth}{M}\setlength{\parskip}{0pt}\setlength{\parindent}{0pt}
  \frenchspacing}{\par\endgroup\smallskip}
\newcommand{\leanline}[2]{\hangindent=\dimexpr#1\leancharwidth+4\leancharwidth\relax
  \hangafter=1\noindent\hspace*{#1\leancharwidth}#2\par}
\newcommand{\leansource}[1]{\par{\footnotesize\color{leanmuted}#1}\par}
\newcommand{\leanlabel}[1]{\par\smallskip\noindent{\small\bfseries #1}\par}
\usepackage[margin=25mm]{geometry}
\usepackage[unicode,hidelinks]{hyperref}
\hypersetup{pdfauthor={LangevinLean Project}}
\setlength{\emergencystretch}{3em}
\setlength{\parskip}{0.4em}
\setlength{\parindent}{0pt}
\allowdisplaybreaks[2]
\raggedbottom
\title{\texorpdfstring{From stationary comparison to the actual invariant-law bounds}{From stationary comparison to the actual invariant-law bounds}}
\author{LangevinLean Project}
\date{}
\begin{document}\maketitle
Draft mathematical exposition; these notes may contain mistakes.


These reusable conclusions identify the comparison target with the actual Gibbs law and its chain with actual Gaussian ULA iterates. The constructed invariant law, its proved attraction and the genuine transport triangle give the exact bias, convergence and accuracy bounds once the concrete comparison is supplied.



Earlier results used here: \cite[Section~3.26]{LangevinLean2026ULASC}.

\section{\texorpdfstring{Actual law identifications and finite transport}{Actual law identifications and finite transport}}


\textbf{Set-up.} Identify the stationary reference law with \(\pi\), the chain law with \(K_h^n\pi\), and the parameter dimension with the physical dimension. The proved stationary comparison then bounds \(W_2(\pi,K_h^n\pi)\) uniformly and proves finiteness. Gaussian ULA iteration preserves probability, and finite second moments imply finite transport from any such iterate to the constructed invariant law \(\pi_h\).



\textbf{Proof.} Substitute the exact law identifications into the comparison coupling. Iteration is the actual Gaussian Markov recursion. Its proved contraction transports the finite initial second-moment distance to the invariant law.

\begin{leanbox}{Lean: \texttt{ULAPoissonComparison}, \texttt{stationary\_\allowbreak{}start\_\allowbreak{}bound}, \texttt{stationary\_\allowbreak{}start\_\allowbreak{}finite} and 2 more}
\leansource{\texttt{ULAPoissonComparison} \quad WassersteinMixing/\allowbreak{}ComparisonTheorems.lean, lines 30–37}
\begin{leancode}
\leanline{0}{structure\ ULAPoissonComparison\ (S\ :\ PoissonProcess\ H\ E\ Q\ P.step\ P.step\_pos)}
\leanline{4}{(score\ :\ H)\ (V\ :\ E\ →\ ℝ)\ extends\ PoissonComparison\ S\ score\ where}
\leanline{2}{dimension\_eq\ :\ P.dimension\ =\ Module.finrank\ ℝ\ E}
\leanline{2}{target\_law\ :\ Q.map\ (reference\ 0)\ =\ ULA.gibbsMeasure\ V}
\leanline{2}{chain\_law\ :\ ∀\ n,\ Q.map\ (chain\ n)\ =}
\leanline{4}{(ULA.ulaLawStep\ V\ P.step)\textasciicircum{}[n]\ (ULA.gibbsMeasure\ V)}
\leanline{0}{\strut}
\leanline{0}{namespace\ ULAPoissonComparison}
\end{leancode}
\leansource{\texttt{stationary\_\allowbreak{}start\_\allowbreak{}bound} \quad WassersteinMixing/\allowbreak{}ComparisonTheorems.lean, lines 43–49}
\begin{leancode}
\leanline{0}{theorem\ stationary\_start\_bound\ (n\ :\ ℕ)\ :}
\leanline{4}{ULA.wasserstein2Real\ (ULA.gibbsMeasure\ V)}
\leanline{6}{((ULA.ulaLawStep\ V\ P.step)\textasciicircum{}[n]\ (ULA.gibbsMeasure\ V))\ ≤\ P.bias}
\end{leancode}
\leansource{\texttt{stationary\_\allowbreak{}start\_\allowbreak{}finite} \quad WassersteinMixing/\allowbreak{}ComparisonTheorems.lean, lines 50–54}
\begin{leancode}
\leanline{0}{theorem\ stationary\_start\_finite\ (n\ :\ ℕ)\ :}
\leanline{4}{ULA.wasserstein2\ (ULA.gibbsMeasure\ V)}
\leanline{6}{((ULA.ulaLawStep\ V\ P.step)\textasciicircum{}[n]\ (ULA.gibbsMeasure\ V))\ ≠\ ⊤}
\end{leancode}
\leansource{\texttt{ula\_\allowbreak{}iterate\_\allowbreak{}isProbability} \quad WassersteinMixing/\allowbreak{}ComparisonTheorems.lean, lines 59–65}
\begin{leancode}
\leanline{0}{theorem\ ula\_iterate\_isProbability\ (P\ :\ Parameters)\ (V\ :\ E\ →\ ℝ)}
\leanline{4}{(hV\ :\ PotentialAssumptions\ P\ V)\ (μ\ :\ Measure\ E)\ [IsProbabilityMeasure\ μ]\ (n\ :\ ℕ)\ :}
\leanline{4}{IsProbabilityMeasure\ ((ULA.ulaLawStep\ V\ P.step)\textasciicircum{}[n]\ μ)}
\end{leancode}
\leansource{\texttt{transport\_\allowbreak{}invariant\_\allowbreak{}finite} \quad WassersteinMixing/\allowbreak{}ComparisonTheorems.lean, lines 69–75}
\begin{leancode}
\leanline{0}{theorem\ transport\_invariant\_finite\ (P\ :\ Parameters)\ (V\ :\ E\ →\ ℝ)}
\leanline{4}{(hV\ :\ PotentialAssumptions\ P\ V)\ (μ\ :\ Measure\ E)\ [IsProbabilityMeasure\ μ]}
\leanline{4}{(hμ\ :\ MemLp\ (id\ :\ E\ →\ E)\ 2\ μ)\ (n\ :\ ℕ)\ :}
\leanline{4}{ULA.wasserstein2\ ((ULA.ulaLawStep\ V\ P.step)\textasciicircum{}[n]\ μ)\ (invariantLaw\ P\ V\ hV)\ ≠\ ⊤}
\end{leancode}
\leanlabel{How the Lean proof follows the argument}
\texttt{ULAPoisson\allowbreak{}Comparison} adds only dimension and actual-law equalities to the generic comparison. The stationary-start theorems substitute these equalities. Probability and finiteness of iterates use the actual Gaussian Markov law and the independently constructed invariant law.
\end{leanbox}


Earlier results used here: \cite[Section~3.17.3]{LangevinLean2026ULASC}; \cite[Section~3.17.4]{LangevinLean2026ULASC}; \cite[Section~3.21.2]{LangevinLean2026ULASC}; \cite[Section~3.26.3]{LangevinLean2026ULASC}; \cite[Section~3.33.11]{LangevinLean2026ULASC}; \cite[Section~3.34.1]{LangevinLean2026ULASC}; \cite[Section~3.36.5]{LangevinLean2026ULASC}; \cite[Section~3.36.7]{LangevinLean2026ULASC}.

\section{\texorpdfstring{The exact invariant bias and finite-time convergence}{The exact invariant bias and finite-time convergence}}


\textbf{Claim.}

\[W_2(\pi,\pi_h)\le6h\sqrt{\kappa\beta d},\qquad W_2(K_h^n\mu,\pi)\le(1-\alpha h)^nW_2(\mu,\pi_h)+6h\sqrt{\kappa\beta d},\quad\kappa=\beta/\alpha.\]

\textbf{Proof.} Apply the transport triangle through \(K_h^n\pi\). Its first leg has the uniform stationary-start bound; its second tends to zero by actual invariant attraction. Taking this proved limit yields the bias. For the finite-time statement, apply the triangle through \(\pi_h\), bound its first leg by the proved ULA contraction and its second by the bias. All laws and finite distances have already been established.

\begin{leanbox}{Lean: \texttt{bias\_\allowbreak{}from\_\allowbreak{}comparison}, \texttt{convergence\_\allowbreak{}from\_\allowbreak{}comparison}}
\leansource{\texttt{bias\_\allowbreak{}from\_\allowbreak{}comparison} \quad WassersteinMixing/\allowbreak{}ComparisonTheorems.lean, lines 81–101}
\begin{leancode}
\leanline{0}{theorem\ bias\_from\_comparison\ (hV\ :\ PotentialAssumptions\ P\ V)\ (C\ :\ ULAPoissonComparison\ S\ score\ V)\ :}
\leanline{4}{ULA.wasserstein2Real\ (ULA.gibbsMeasure\ V)\ (invariantLaw\ P\ V\ hV)\ ≤}
\leanline{6}{6\ *\ P.step\ *\ Real.sqrt\ (P.condition\ *\ P.beta\ *\ P.dimension)}
\end{leancode}
\leansource{\texttt{convergence\_\allowbreak{}from\_\allowbreak{}comparison} \quad WassersteinMixing/\allowbreak{}ComparisonTheorems.lean, lines 105–124}
\begin{leancode}
\leanline{0}{theorem\ convergence\_from\_comparison\ (hV\ :\ PotentialAssumptions\ P\ V)}
\leanline{4}{(C\ :\ ULAPoissonComparison\ S\ score\ V)\ (μ\ :\ Measure\ E)\ [IsProbabilityMeasure\ μ]}
\leanline{4}{(hμ\ :\ MemLp\ (id\ :\ E\ →\ E)\ 2\ μ)\ (n\ :\ ℕ)\ :}
\leanline{4}{ULA.wasserstein2Real\ ((ULA.ulaLawStep\ V\ P.step)\textasciicircum{}[n]\ μ)\ (ULA.gibbsMeasure\ V)\ ≤}
\leanline{6}{P.contraction\ \textasciicircum{}\ n\ *\ ULA.wasserstein2Real\ μ\ (invariantLaw\ P\ V\ hV)\ +}
\leanline{8}{6\ *\ P.step\ *\ Real.sqrt\ (P.condition\ *\ P.beta\ *\ P.dimension)}
\end{leancode}
\leanlabel{How the Lean proof follows the argument}
\texttt{bias\_\allowbreak{}from\_\allowbreak{}comparison} passes the actual invariant-attraction limit through a uniform triangle bound. \texttt{convergence\_\allowbreak{}from\_\allowbreak{}comparison} combines \texttt{invariant\allowbreak{}Law\_\allowbreak{}mixing}, the proved transport triangle and that bias. Invariant existence, contraction and limit convergence are not comparison fields.
\end{leanbox}


Earlier results used here: \cite[Section~3.17.4]{LangevinLean2026ULASC}; \cite[Section~3.21.2]{LangevinLean2026ULASC}; \cite[Section~3.26.13]{LangevinLean2026ULASC}; \cite[Section~3.26.15]{LangevinLean2026ULASC}; \cite[Section~3.26.3]{LangevinLean2026ULASC}; \cite[Section~3.36.7]{LangevinLean2026ULASC}.

\section{\texorpdfstring{The source accuracy scale}{The source accuracy scale}}


\textbf{Claim.} With the source step choice and its logarithmic iteration threshold, \(\sqrt\alpha\,W_2(K_h^n\mu,\pi)\le\varepsilon\).



\textbf{Proof.} Multiply the actual convergence bound by \(\sqrt\alpha\). The step choice makes the bias contribution at most \(\varepsilon/2\); the iteration threshold makes the geometric transient at most \(\varepsilon/2\). These exact numerical estimates are proved independently in the complexity calculation. Positive step together with the displayed step formula excludes the zero-dimensional exceptional branch when this accuracy statement is used.

\begin{leanbox}{Lean: \texttt{accuracy\_\allowbreak{}from\_\allowbreak{}comparison}}
\leansource{\texttt{accuracy\_\allowbreak{}from\_\allowbreak{}comparison} \quad WassersteinMixing/\allowbreak{}ComparisonTheorems.lean, lines 128–150}
\begin{leancode}
\leanline{0}{theorem\ accuracy\_from\_comparison\ (hV\ :\ PotentialAssumptions\ P\ V)}
\leanline{4}{(C\ :\ ULAPoissonComparison\ S\ score\ V)\ (μ\ :\ Measure\ E)\ [IsProbabilityMeasure\ μ]}
\leanline{4}{(hμ\ :\ MemLp\ (id\ :\ E\ →\ E)\ 2\ μ)\ (ε\ :\ ℝ)\ (hε\ :\ 0\ <\ ε)}
\leanline{4}{(hstep\ :\ P.step\ =\ paperStep\ P.beta\ P.dimension\ ε)\ (n\ :\ ℕ)}
\leanline{4}{(hn\ :\ iterationPrefactor\ P.alpha\ P.beta\ P.dimension\ ε\ *}
\leanline{6}{Real.log\ (1\ +\ 2\ *\ Real.sqrt\ P.alpha\ *}
\leanline{8}{ULA.wasserstein2Real\ μ\ (invariantLaw\ P\ V\ hV)\ /\ ε)\ ≤\ (n\ :\ ℝ))\ :}
\leanline{4}{Real.sqrt\ P.alpha\ *}
\leanline{6}{ULA.wasserstein2Real\ ((ULA.ulaLawStep\ V\ P.step)\textasciicircum{}[n]\ μ)\ (ULA.gibbsMeasure\ V)\ ≤\ ε}
\end{leancode}
\leanlabel{How the Lean proof follows the argument}
\texttt{accuracy\_\allowbreak{}from\_\allowbreak{}comparison} applies the exact numerical theorem \texttt{paper\_\allowbreak{}complexity\_\allowbreak{}bound} to the actual initial distance from the constructed invariant law. The proof derives the needed positive dimension from the positive step and its prescribed formula.
\end{leanbox}


Earlier results used here: \cite[Section~3.17.4]{LangevinLean2026ULASC}; \cite[Section~3.26.3]{LangevinLean2026ULASC}; \cite[Section~3.36.7]{LangevinLean2026ULASC}.

\begin{thebibliography}{ULA-LSI}
\bibitem[ULA-SC]{LangevinLean2026ULASC}
LangevinLean Project. \emph{Unadjusted Langevin under Strong Convexity}. LL-01, 2026. Draft report, content revision d34ead06987c. \url{https://odakyildiz.github.io/LangevinLean/docs/pdf/UnadjustedLangevin/FullProof.pdf}.
\end{thebibliography}

\end{document}
